\section{Research Questions}
\label{sec:rq}

We organize the study around six empirical questions.
They are answered descriptively from the Phase~2B evidence pack; none is treated as rhetorically settled in advance.

\begin{description}[style=nextline,leftmargin=1.2cm,font=\normalfont\bfseries]
  \item[RQ1] Does authentic ISV corpus priming produce a measurable shift in Polish-to-ISV translation outputs under the canonical coverage metric?
  \item[RQ2] Does this shift persist when the source text has lower thematic overlap with the priming corpus (LOW)?
  \item[RQ3] Does a positive paired shift persist for a domain-unseen, technical/expository source (UNSEEN)?
  \item[RQ4] How heterogeneous are priming effects across the seven model configurations?
  \item[RQ5] What qualitative textual changes accompany positive or negative shifts in canonical coverage?
  \item[RQ6] Which observed changes appear related to corpus-opening copying, lexical adoption, or orthographic normalization?
\end{description}

RQ1--RQ3 address magnitude and direction of $\Delta$ across regimes.
RQ4 addresses configuration dependence.
RQ5--RQ6 connect metric movement to audited textual behaviour without asserting causal mechanisms.


The RQs are ordered from measurement (RQ1) through generalization across sources (RQ2--RQ3) to heterogeneity and mechanism-facing description (RQ4--RQ6).
This order mirrors the analysis workflow actually used in the repository: regime aggregates first, then combined synthesis, then qualitative audits.
